\documentclass[11pt,a4paper]{article}
\usepackage[utf8]{inputenc}
\usepackage[T1]{fontenc}
\usepackage[left=25mm,right=25mm,top=25mm,bottom=25mm]{geometry}
\usepackage{graphicx}
\usepackage{caption}
\usepackage{amsmath}
\usepackage[hidelinks]{hyperref}
\usepackage{booktabs}
\captionsetup{font=small}
\hypersetup{pageanchor=false}
\begin{document}


% Cover page
 \begin{titlepage}
    \begin{center}
    {\includegraphics[width=1\textwidth]{logo.png}\par}
       \vspace{0.5 cm}
       
    {\bfseries\LARGE Universidad de Zaragoza  \par}
    \vspace{0.5 cm}
    
    {\scshape\Large Facultad de Ingenier\'ia\par}
    \vspace{0.5 cm}
    
    {\scshape\LARGE  Artificial Intelligence - Lab Session 3: Neural Networks and CNNs\par}
    \vspace{1.5 cm}

    \vfill
    {\Large Jorge Ruiz González (826685) \par}
    {\Large Hugo López Sola (875455) \par}
    
    \vfill
    
    {\Large October 14th, 2026 \par}
    \end{center}
    \end{titlepage}\hypersetup{pageanchor=true}

    \tableofcontents


\clearpage
\section{Aim and evaluation}
In Task 1, we first train the two-layer network on the synthetic toy data, then
on CIFAR-10. For CIFAR-10 we use 49,000 training images, 1,000 validation images
and 10,000 test images. In Tasks 2 and 3, we train CNNs on five food classes:
cheesecake, chicken quesadilla, chicken wings, churros and donuts. There are
1,000 images per class, split into 4,000 training and 1,000 validation images
with seed 123.

We use validation accuracy to compare the models. The food dataset has no
separate test set. Its results cannot be directly compared with CIFAR-10,
since the images and classes are different.

\clearpage
\section{Task 1: Two-layer neural network}
The network applies a fully connected layer, ReLU and a second fully connected
layer. The last layer produces class scores; softmax cross-entropy measures the
classification error. L2 regularization penalizes large weights, and stochastic
gradient descent updates the parameters using minibatches.

On the five-sample, three-class toy problem, training reaches a loss of 0.0367,
below the requested 0.2. The toy initialization uses a standard deviation of
0.1; CIFAR-10 retains the starter's default initialization. CIFAR images are
flattened into 3,072 features and centred using only the training mean.

The initial CIFAR model uses 50 hidden units and 1,000 iterations, obtaining
28.60\% validation accuracy. We then run 20 random trials of 2,000 iterations,
using learning rates $\{10^{-4},5\times10^{-4},10^{-3}\}$, hidden sizes
$\{50,100,150\}$ and regularization strengths $\{0.25,0.5,0.75\}$.
We leave out $10^{-2}$ because it caused numerical overflow in the earlier trials.

\begin{table}[ht]
\centering
\small
\caption{Selected two-layer network on CIFAR-10.}
\begin{tabular}{lrrr}
\toprule
Configuration & Training & Validation & Test\\
\midrule
100 hidden units, $lr=10^{-3}$, $reg=0.25$ & 54.52\% & 52.50\% & 49.96\%\\
\bottomrule
\end{tabular}
\end{table}

We evaluate the selected model on the test set only after the search. To reduce
overfitting, we would try more training data or stronger L2 regularization.
Adding hidden units can make overfitting worse. A smaller learning rate changes
the training updates but does not directly solve this problem.

\clearpage
\section{Task 2: Simple CNN}
The baseline resizes food images to $150\times150$, rescales their intensities
and uses three convolutional/pooling blocks with 16, 32 and 64 filters, followed
by a 128-unit dense layer and five output scores. We compare it with dropout,
then add horizontal flips, small rotations, zoom and contrast changes. The
larger model uses 32/64/128 filters and 256 dense units; all kernels remain
$3\times3$.

\begin{table}[ht]
\centering
\small
\caption{Final-epoch results saved in the main Task 2 notebook.}
\begin{tabular}{lrrr}
\toprule
Model & Epochs & Training & Validation\\
\midrule
Baseline & 10 & 99.10\% & 41.30\%\\
Dropout & 10 & 72.35\% & 48.80\%\\
Dropout + augmentation & 20 & 52.93\% & 53.60\%\\
Larger CNN + augmentation & 20 & 56.15\% & 54.40\%\\
\bottomrule
\end{tabular}
\end{table}

Without dropout, training accuracy is almost 100\%, but validation accuracy is
only 41.30\%. This shows clear overfitting. Dropout reduces the gap and raises
validation accuracy to 48.80\%.

With augmentation, validation accuracy is 47.60\% at epoch 10 and 53.60\% at
epoch 20. We therefore cannot attribute the whole improvement to augmentation:
this model also trains for longer. Its training and validation accuracies are
similar, but training uses dropout and augmented images while validation does
not. This alone is not enough to conclude that the model underfits.

The separate longer-run notebook reaches 57.40\% final validation accuracy with
the larger CNN after 40 epochs (61.30\% at its best epoch). Since this is a
separate run, the difference is not a controlled comparison of epoch counts.

\clearpage
\section{Task 3: Transfer learning and fine-tuning}
We load MobileNetV2 with ImageNet weights and remove its original classifier.
Images are resized to $224\times224$ and receive the model's matching
preprocessing. Global average pooling, dropout and a five-output dense layer
form the new classifier. The output is logits, so sparse categorical
cross-entropy uses \texttt{from\_logits=True}.

First, the backbone is frozen and the new classifier is trained with Adam at
$10^{-3}$. We then unfreeze the backbone and recompile at $10^{-4}$. Two
variations restart from the same saved frozen-head weights: full fine-tuning at
$10^{-5}$, and fine-tuning only the last 20 backbone layers at $10^{-5}$, with
BatchNormalization layers frozen. Backbone BatchNormalization statistics remain
in inference mode throughout, as in the starter.

We run five epochs per stage on CPU, with seed 123 and batch size 4. Both
variations use new optimizers and callbacks. We save the weights with the
lowest validation loss and set early-stopping patience to five epochs.
Table~\ref{tab:transfer} shows the final epoch and the selected checkpoint
separately, since they are not always the same.

\begin{table}[ht]
\centering
\small
\caption{Task 3 results. Selected checkpoints minimize validation loss.}
\label{tab:transfer}
\begin{tabular}{lrrrr}
\toprule
Stage, learning rate & Epochs & Final train & Final val. & Selected val.\\
\midrule
Frozen backbone, $10^{-3}$ & 5 & 89.92\% & 90.90\% & --\\
Full fine-tuning, $10^{-4}$ & 5 & 90.50\% & 71.70\% & --\\
Full fine-tuning, $10^{-5}$ & 5 & 82.13\% & 79.10\% & 88.30\%\\
Last 20 layers, $10^{-5}$ & 5 & 95.95\% & 93.10\% & 93.10\%\\
\bottomrule
\end{tabular}
\end{table}

The frozen backbone already gives 90.90\% validation accuracy. Unfreezing all
layers at $10^{-4}$ makes the result worse: it falls to 71.70\%, even though
training accuracy reaches 90.50\%. Lowering the learning rate to $10^{-5}$
gives 88.30\% at the selected checkpoint, but only 79.10\% at the final epoch.

Training only the last layers works best here, reaching 93.10\%. These results
do not support unfreezing every layer by default. We would keep the partial
fine-tuning model for this dataset. The comparison uses one split and five
epochs per stage, so it does not show which strategy would always be best.

The ImageNet weights are necessary for this transfer-learning experiment.
Loading the architecture with random weights would mean training from scratch.
Compared with the simple CNN's 54.40\%, the pretrained model gives much higher
validation accuracy, although the image resolution and training budgets differ.
It also requires the correct preprocessing and a suitable fine-tuning learning
rate.

\end{document}
